\reading{LTR-8}{The Failure Mechanics of Dealer Banks}{DEFER}{0}

\begin{lrkeybox}[Learning objectives — official 2026 spine wording]
\begin{enumerate}[label=\textbf{\alph*.}]
\item Compare and contrast the major lines of business in which dealer banks operate and the risk factors they face in each line of business.
\item Identify situations that can cause a liquidity crisis at a dealer bank and explain responses that can mitigate these risks.
\item Assess policy measures that can alleviate firm-specific and systemic risks related to large dealer banks.
\end{enumerate}
{\footnotesize\color{FRMGrey}Source: Darrell Duffie, ``The Failure Mechanics of Dealer Banks,'' \textit{Journal of Economic Perspectives} 24:1, 2010.}
\end{lrkeybox}

\begin{lrtrapbox}[Label collision]
The Lecture layer labels this reading \texttt{LO 66.a--c} --- the same \texttt{66} prefix the Crash layer uses for \textbf{LTR-3 Early Warning Indicators}. Renumbered to the spine.
\end{lrtrapbox}

\subsection*{Basics: what a dealer bank is}

\begin{itemize}
\item Dealer banks are typically parts of \textbf{large financial organizations} that operate other financial businesses. Many large dealer banks have conventional commercial banking operations, including deposit taking as well as lending to corporations and consumers.
\item The major functions of large dealer banks are: intermediation of the markets for securities, securities lending, repurchase agreements and derivatives; \textbf{prime brokerage} for hedge funds; \textbf{asset management} for institutional and wealthy individual investors; and \textbf{proprietary trading} --- speculation on their own accounts.
\end{itemize}

% =====================================================================
\lo{LTR-8 a}{Compare and contrast the major lines of business in which dealer banks operate and the risk factors they face in each line of business.}

\begin{center}
\scriptsize
\begin{longtable}{@{}p{2.9cm} p{5.4cm} p{5.2cm}@{}}
\toprule
\rowcolor{FRMSoft}\textbf{Line of business} & \textbf{Function} & \textbf{Risk factors} \\
\midrule
\endfirsthead
\toprule
\rowcolor{FRMSoft}\textbf{Line of business} & \textbf{Function} & \textbf{Risk factors} \\
\midrule
\endhead
\textbf{OTC derivatives} &
Act as counterparties; transfer risk among counterparties through new contracts --- interest rate swaps, CDOs, CMOs, CDSs. &
Counterparty default, which raises systemic risk; high exposure to \emph{other dealer banks}. \\
\addlinespace
\textbf{Short-term repo market} &
Finance assets with short-term loans. High leverage due to the lack of capital requirements pre-crisis, with subprime mortgages used as collateral. &
Counterparty questions solvency and does not renew repos; repo creditors sell the collateral, producing a liquidity crisis. \\
\addlinespace
\textbf{Investment banking} &
Underwrite securities issuance; M\&A advisory; merchant banking; commodity trading. &
Issuers take business elsewhere if solvency is questioned, producing liquidity strain; reduced cash inflows; disrupted liquidity of existing securities. \\
\addlinespace
\textbf{Prime brokerage} &
Custody, clearing, and securities lending for hedge funds. &
Hedge funds demand \textbf{cash} margin loans, and the dealer can no longer use their collateral. Flight of clients reduces the liquidity pool. Systemic risk if hedge funds are not diversified across prime brokers. \\
\addlinespace
\textbf{Off-balance-sheet entities} &
Internal hedge funds, private equity partnerships, and SPEs; before the crisis, sponsored SIVs financing mortgages with short-term debt. &
Losses in these entities strain the dealer bank's finances. \textbf{Reputation risk} if the dealer bank does not support them --- which is why dealer banks supported SIVs even though the exposure was off balance sheet. \\
\addlinespace
\textbf{Traditional banking functions} &
Deposit gathering; corporate and consumer lending. &
Similar risks to traditional banks: deposit run, credit risk. Before 2007--2009, dealer banks \textbf{lacked access to central bank borrowing and federal deposit insurance}. \\
\bottomrule
\end{longtable}
\end{center}
\figcap{Read the risk column downward and the same word appears in five of six rows: \emph{solvency}. Every line of business fails through the same channel --- a counterparty, client or depositor doubting solvency and withdrawing. What differs is who withdraws and what they take with them.}

\begin{lrkeybox}[Economies and diseconomies of scope]
Large dealer banks offer a variety of services, which creates \textbf{economies of scope} in areas such as technology, innovation and marketing. But the 2007--2009 crisis revealed \textbf{diseconomies of scope} in risk management and corporate governance: the complexity hindered risk management, and executive management and boards often lacked a complete understanding of and control over the risks their organizations took.

\textbf{Examples.} Bear Stearns and Lehman Brothers relied heavily on overnight repos with high leverage ratios, holding these assets on their balance sheets without sufficient capital. Inadequate risk management contributed to their insolvency when the value of their assets was questioned. Improved awareness and more appropriate risk models could potentially have prevented it.
\end{lrkeybox}

\begin{lrtrapbox}[Economies of scope are real; the diseconomies are in governance]
The reading does \emph{not} say the breadth of services was itself the problem. It says the benefit is in technology and marketing and the cost is in \textbf{risk management and corporate governance}. An option asserting that dealer banks gain no economies of scope, or that the diseconomies were operational rather than governance-related, has corrupted one half of a two-sided claim.
\end{lrtrapbox}

\begin{lrkeybox}[Where dealer banks sit relative to the safety net]
Large dealer banks operate \textbf{outside traditional bank failure resolution mechanisms} and provide their services under holding companies. Some engage in off-balance-sheet financing by selling loans to \term{special purpose entities (SPEs)}, compensated with debt issued to third-party investors; the SPEs use cash flows from those assets to pay principal and interest. Before the crisis, dealer banks sponsored \term{structured investment vehicles (SIVs)} --- a form of SPE financing mortgages and other debt with short-term debt. The crisis revealed the risk: off-balance-sheet assets and debt obligations were not accounted for, leading to liquidity and solvency problems, and the dealer banks had to support the SIVs to protect their reputation.
\end{lrkeybox}

% =====================================================================
\lo{LTR-8 b}{Identify situations that can cause a liquidity crisis at a dealer bank and explain responses that can mitigate these risks.}

\subsection{The four liquidity concerns, and what mitigates each}

\begin{center}
\small
\begin{tabularx}{\textwidth}{@{}p{3.0cm} X p{4.0cm}@{}}
\toprule
\rowcolor{FRMSoft}\textbf{Concern} & \textbf{The mechanism} & \textbf{Mitigant} \\
\midrule
\textbf{1. Counterparty risk} &
When counterparties doubt a dealer bank's solvency they reduce exposure --- exiting derivatives contracts, requesting \textbf{novations}, restructuring OTC derivatives, or demanding cash --- worsening the liquidity crisis. &
\textbf{Central clearing}, effective for \emph{standardized} derivatives. \\
\addlinespace
\textbf{2. Repo market risk} &
If creditors lose trust they may not renew short-term funding agreements, forcing the dealer bank to sell assets quickly, potentially at a loss, and potentially triggering immediate collateral sales by repo counterparties. &
\textbf{Laddering maturities of liabilities}, so a significant portion of debt does not need refinancing overnight. \\
\addlinespace
\textbf{3. Prime brokerage risk} &
If prime brokerage clients withdraw their cash and securities, the dealer bank's liquidity pool shrinks. Hedge funds may request cash margin loans backed by securities. The exit of clients reduces collateral, potentially causing a \textbf{systemic collateral shortage}. &
\textbf{Diversifying} prime brokerage services, which lessens client concentration risk. Hedge funds may also diversify their prime brokers. \\
\addlinespace
\textbf{4. Clearing bank risk} &
Clearing banks may stop extending credit, or refuse to process transactions insufficiently funded by a dealer bank's cash account, if solvency is questioned. &
\textbf{``Emergency bank'' proposals}; regulatory support such as the Primary Dealer Credit Facility. \\
\bottomrule
\end{tabularx}
\end{center}

\begin{lrtrapbox}[Match the mitigant to the concern, not to the general theme]
All four mitigants sound plausible against all four concerns, which is exactly why they are offered crosswise. The pairings the reading makes are: counterparty $\rightarrow$ \textbf{central clearing}; repo $\rightarrow$ \textbf{laddering liabilities}; prime brokerage $\rightarrow$ \textbf{diversification}; clearing bank $\rightarrow$ \textbf{emergency bank}. Also note the qualifier on the first: central clearing works for \textbf{standardized} contracts, and an option dropping that word overstates the fix.
\end{lrtrapbox}

\begin{lrkeybox}[Comparison with a retail bank run]
\begin{itemize}
\item \textbf{Economic principles.} Liquidity crises at dealer banks share the same driver as traditional retail bank runs: concerns about solvency.
\item \textbf{Institutional differences.} The institutional mechanisms and the systemic impacts differ significantly between dealer banks and retail banks.
\item \textbf{Unique pressures.} Dealer banks face liquidity pressure from \emph{multiple functions at once} --- OTC derivatives, repo, prime brokerage, clearing --- creating challenges that traditional banks do not encounter. Highly leveraged dealer banks are especially susceptible when repo counterparties refuse to renew positions.
\end{itemize}
\end{lrkeybox}

\begin{lrtrapbox}[Same cause, different plumbing]
The standard four-statement build on this point keeps ``driven by solvency concerns'' (true) and flips ``institutional mechanisms are similar'' (false), or vice versa. The reading's position is: \textbf{same economics, different institutions, larger systemic impact}.
\end{lrtrapbox}

% =====================================================================
\lo{LTR-8 c}{Assess policy measures that can alleviate firm-specific and systemic risks related to large dealer banks.}

\begin{center}
\small
\begin{tabularx}{\textwidth}{@{}p{3.8cm} X@{}}
\toprule
\rowcolor{FRMSoft}\textbf{Measure} & \textbf{What it does} \\
\midrule
\textbf{Public Private Investment Partnership (PPIP), 2009} &
Created to help dealer banks recover from the financial crisis. Aimed to address the problem of \textbf{``toxic'' assets} by offering favourable financing terms and \textbf{absorbing losses for buyers}. \\
\addlinespace
\textbf{Central bank secured lending facilities} &
During the crisis, the \textbf{Federal Reserve and the Bank of England} provided secured lending facilities to dealer banks when they could not get credit elsewhere, such as in the repo market. \\
\addlinespace
\textbf{Tri-party repo utility} &
A proposed solution to manage repo positions more effectively, operating with \textbf{stricter standards}. \\
\addlinespace
\textbf{``Emergency bank''} &
A proposal to establish an entity to oversee the \textbf{unwinding of dealer banks' repo positions} during times of liquidity difficulty, with central bank support. \\
\addlinespace
\textbf{Higher capital requirements} &
Expected to rise, \textbf{including off-balance-sheet positions}, to reduce dealer banks' leverage. \\
\addlinespace
\textbf{Central clearing} &
Suggested to minimize the risk of OTC derivatives counterparties leaving questionable dealer banks. \\
\addlinespace
\textbf{Bridge banks} &
Some large dealer banks were considered \textbf{``too big to fail''}; the proposal included providing bridge banks, similar to those used for traditional banks, to address solvency concerns and systemic risks. \\
\bottomrule
\end{tabularx}
\end{center}
\figcap{Two of these were actually done (PPIP, central bank facilities) and five were \emph{proposals} (tri-party utility, emergency bank, higher capital, central clearing, bridge banks). GARP's favourite corruption on a policy list is to promote a proposal into an implemented measure, or to attribute an implemented one to the wrong authority.}

\begin{lrtrapbox}[Two attribution details worth locking]
\begin{itemize}
\item PPIP absorbed losses \textbf{for buyers} of toxic assets --- not for the dealer banks selling them.
\item The secured lending facilities were provided by the \textbf{Federal Reserve and the Bank of England}. Naming the ECB, or the Treasury, is the role-misassignment build.
\end{itemize}
\end{lrtrapbox}

\begin{lrtrapbox}[Consolidated trap summary — LTR-8]
\begin{itemize}
\item \textbf{Role.} Five of six business lines fail through a \emph{solvency doubt} channel; what differs is who withdraws. Match the risk factor to the line, not to the theme.
\item \textbf{Sibling.} Counterparty $\rightarrow$ central clearing; repo $\rightarrow$ laddering liabilities; prime brokerage $\rightarrow$ diversification; clearing bank $\rightarrow$ emergency bank.
\item \textbf{Scope.} Central clearing mitigates counterparty risk only for \textbf{standardized} derivatives.
\item \textbf{Scope.} Economies of scope are in technology and marketing; \textbf{dis}economies are in risk management and corporate governance.
\item \textbf{Scope.} Before 2007--2009 dealer banks \emph{lacked} access to central bank borrowing and federal deposit insurance, and operate outside traditional failure resolution mechanisms.
\item \textbf{Role.} PPIP absorbed losses for \emph{buyers}. Secured lending came from the \emph{Fed and the Bank of England}.
\item \textbf{Sibling.} Same economics as a retail bank run, but different institutional mechanisms and systemic impact.
\item \textbf{Scope.} Dealer banks supported their SIVs for \textbf{reputation}, despite the exposure being off balance sheet --- the risk was never genuinely transferred.
\end{itemize}
\end{lrtrapbox}
